Dear Editor,

We submit our manuscript entitled "World Model for Brain-Blood Mapping and Virtual Knockout Validation Uncovers Drug Repurposing Opportunities and Diagnostic Biomarker Candidates in Alzheimer's Disease" for consideration as an Article in *Nature Machine Intelligence*.

**Conceptual advance.** A fundamental limitation of AI in biomedical research is that most models serve as black-box predictors — they map inputs to outputs but do not yield differentiable, interpretable representations of the underlying biological dynamics. This prevents in silico perturbation experiments: the ability to ask "if I perturb gene *i* in the brain, how does gene *j* in the blood change?" We introduce a *molecular world model*, inspired by the world model paradigm from reinforcement learning, whose core innovation is differentiability. Using a Neural Ordinary Differential Equation as its engine and Sinkhorn Optimal Transport for distribution matching, the world model learns a continuous transformation between brain and blood molecular compartments. Because this transformation is differentiable, its Jacobian sensitivity matrix can be computed via automatic differentiation — yielding, for every gene pair, a quantitative prediction of cross-tissue perturbation response. This Jacobian is the computational analog of a perturbation screen and is a capability that no discrete OT, linear regression, or deep learning classifier provides.

**New capability, not benchmark improvement.** The Jacobian sensitivity network achieves concordance with independent functional knockout validation at AUC = 0.893, compared to 0.560 for the best linear alternative. This is not a marginal gain on a standard benchmark; it is a qualitatively new capability — the differentiable world model is the only approach that produces a biologically interpretable sensitivity matrix from distribution-level training data alone. We further demonstrate three world-model-exclusive properties: continuous trajectory interpolation ($t \in [0,1]$), reversibility (round-trip Pearson $r = 0.974$), and cross-species transferability (mouse-trained model applied to 745 human ADNI subjects, 1,961 of 2,000 genes matched).

**Scientific discovery enabled by AI.** Applied to Alzheimer's disease, the world model identifies three translational outputs validated in the ADNI clinical cohort ($n > 600$):

1. *MAPT* as a dual-qualified diagnostic biomarker (single-gene AUC = 0.719; Cox HR = 1.14, $p = 0.003$), confirmed by Kaplan–Meier survival analysis. Aging pseudotime analysis further shows MAPT expression declining across the lifespan, consistent with cumulative neuronal loss.

2. Three therapeutic targets with clinical-stage drugs: CSF3R (Phase II filgrastim), MMP9 (Phase III andecaliximab), and CXCR2 (Phase III danirixin) — each passing multi-tier virtual knockout + survival + drug-evidence validation.

3. Multi-omics complementarity: the proteomics layer captures the tau pathology diagnostic axis, while the transcriptomics layer captures the immune/inflammation therapeutic axis — a complementary architecture with direct implications for clinical trial design.

**Why *Nature Machine Intelligence*.** This work demonstrates a principled AI method — a differentiable world model that enables a new class of in silico perturbation experiments — applied to a major biomedical problem (AD drug and biomarker discovery) with rigorous benchmarking, cross-species validation, and human clinical evidence. The framework is designed as a reusable computational paradigm: the world model's learned transformation captures evolutionarily conserved regulatory structure and can be retrained on new cross-tissue datasets without architectural modification, supporting the broader vision of AI-driven virtual cell discovery.

All data are from public repositories (GEO, ADNI, BioStudies, figshare). Code will be deposited on GitHub upon acceptance. We have no competing interests to declare. We confirm that this manuscript has not been published and is not under consideration elsewhere.

Thank you for your consideration.

Sincerely,

Lianghua Wang (corresponding author)
Department of Biochemistry and Molecular Biology
Naval Medical University, Shanghai 200433, China
Email: lhwang@smmu.edu.cn
